\documentclass{article} 
\input{../../lib.sty} 
 
\title{\texttt{fn conservative\_continuous\_gaussian\_tail\_to\_alpha}} 
\author{\MichaelShoemateAuthor} 
 
\begin{document} 
\maketitle 
 
Proof for \rustdoc{accuracy/tail\_bounds/fn}{conservative\_continuous\_gaussian\_tail\_to\_alpha}. 
 
 
\section{Hoare Triple} 
\subsection*{Precondition} 
\subsubsection*{Compiler-verified} 
\begin{itemize} 
    \item Argument \texttt{scale} is of type \texttt{f64} 
    \item Argument \texttt{tail} is of type \texttt{f64} 
\end{itemize} 
 
\subsubsection*{Caller-verified} 
\begin{itemize} 
    \item $\texttt{scale} > 0$ 
    \item $\texttt{tail} > 0$ 
\end{itemize} 
 
 
\subsection*{Pseudocode} 
\lstinputlisting[language=Python,firstline=2,escapechar=|]{./pseudocode/conservative_continuous_gaussian_tail_to_alpha.py} 
 
\subsection*{Postcondition} 
 
\begin{theorem}
    Returns \texttt{Ok(out)}, where \texttt{out} is at least $\Pr[X > \texttt{tail}]$ 
    for $X \sim \mathcal{N}(0, \texttt{scale})$, assuming $\texttt{tail} > 0$, 
    or \texttt{Err(e)} if any numerical computation overflows. 
\end{theorem} 
 
\begin{definition} 
    \label{gaussian} 
    Define $X \sim \mathcal{N}(0, \texttt{scale})$, where $\texttt{scale} > 0$. Then $X$ follows the continuous gaussian distribution: 
    \begin{equation} 
        f(x) = \frac{1}{\texttt{scale} \sqrt{2 \pi}} e^{-\frac{1}{2}\left( \frac{x}{\texttt{scale}}\right)^2} 
    \end{equation} 
\end{definition} 
 
\begin{definition} 
    \label{erf} 
    The error function is defined as: 
    \begin{equation} 
        \mathrm{erf}(z) = \frac{2}{\sqrt{\pi}} \int_{0}^{z} e^{-t^2} dt 
    \end{equation} 
\end{definition} 
 
\begin{definition} 
    \label{erfc} 
    The complementary error function is defined as: 
    \begin{equation} 
        \mathrm{erfc}(z) = 1 - \mathrm{erf}(z) 
    \end{equation} 
\end{definition} 
 
\begin{lemma} 
    \label{erfc-err} 
    The implementation of \texttt{erfc} differs from a conservatively rounded implementation by no greater than one 32-bit float ulp. 
\end{lemma} 
 
\begin{proof} 
    The following code conducts an exhaustive search. 
    % the adjacent erfc_err_analysis.py file is checked-in, and runnable 
    \label{sec:erfc-err-check} 
    \lstinputlisting[language=Python,firstline=2]{./erfc_err_analysis.py} 
 
    Upon completion, the greatest discovered error is at most 1 ulp. 
\end{proof} 
 
\begin{theorem} 
    Assume $X \sim \mathcal{N}(0, \texttt{scale})$ and $\texttt{tail} > 0$. Then 
    \begin{equation} 
        \alpha = P[X > \texttt{tail}] = \frac{1}{2} \mathrm{erfc}\left(\frac{\texttt{tail}}{\texttt{scale} \sqrt{2}}\right) 
    \end{equation} 
\end{theorem} 
 
\begin{proof} 
    The strict postcondition $P[X > \texttt{tail}]$ is intentional, because dependent proofs use this exact event.
    The argument in this theorem is about the ideal continuous gaussian random variable $X$.
    For this reference distribution, continuity implies $P[X > \texttt{tail}] = P[X \ge \texttt{tail}]$.
    This equality is used only within the theorem, and not for any discretized or postprocessed implementation,
    where $P[X > \texttt{tail}]$ and $P[X \ge \texttt{tail}]$ may differ due to point masses.

    \begin{align*} 
        \alpha &= P[X > \texttt{tail}] \\ 
        &= P[X \ge \texttt{tail}] && \text{because } X \text{ is continuous} \\ 
        &= \int_{\texttt{tail}}^{\infty} \frac{1}{\texttt{scale} \sqrt{2 \pi}} e^{-\frac{1}{2}\left( \frac{x}{\texttt{scale}}\right)^2} dx && \text{by } \ref{gaussian} \\ 
        &= \frac{1}{\sqrt{\pi}} \int_{\texttt{tail} / (\texttt{scale} \sqrt{2})}^{\infty} e^{-u^2} du && \text{with } u = \frac{x}{\texttt{scale} \sqrt{2}} \\ 
        &= \frac{1}{2} \left(1 - \mathrm{erf}\frac{\texttt{tail}}{\texttt{scale} \sqrt{2}}\right) && \text{by } \ref{erf} \\ 
        &= \frac{1}{2} \mathrm{erfc}\left(\frac{\texttt{tail}}{\texttt{scale} \sqrt{2}}\right) && \text{by } \ref{erfc} 
    \end{align*} 
 
    The implementation computes $\frac{\texttt{tail}}{\texttt{scale} \sqrt{2}}$ with conservative downward rounding, so the input to \texttt{erfc} is no greater than the exact value. 
    Since \texttt{erfc} is monotonically decreasing, this yields a value no smaller than the exact tail probability expression above. 
    The outcome of \texttt{erfc} is then increased by one 32-bit float ulp, which compensates for implementation error by \ref{erfc-err}. 
    Therefore the returned result is at least $\Pr[X > \texttt{tail}]$. 
\end{proof} 
 
\end{document} 
